\subsection{Extension by Finite Axis Search}
\label{app:axis-extensions}
The following result extends semantic coverage using finite search over
axis choices, rather than the direct constructor of
Theorem~\ref{thm:complete-synthesis}.

\begin{theorem}[Extended axis synthesis]
\label{thm:extended-axis-synthesis}
Under Section~\ref{sec:formal-scope}'s input and rewrite assumptions,
allow axis reversal, broadcast inserting axes of supplied extents at
least two, and fixed nonzero rational unary scaling.
If legal axis choices are finitely enumerable from metadata and child
shapes while preserving supplied fields, a terminating metadata-only
search finds a valid RHS completion if and only if one exists.
\end{theorem}

\begin{proof}
Represent each component by $N^+(T)=(s,\Pi,w)$: its ordered output
shape, port partition with canonical relative directions, and nonzero
rational weight. Broadcast introduces independent axes; scaling
multiplies $w$. Summed classes containing no input or output port
contribute their extents to $w$ and are removed.

This representation is exact. Equal representations give equal
computations. Different shapes are inequivalent; different signed
partitions are distinguished by one-hot inputs using endpoint
coordinates $0,n-1$, since all extents are at least two. For identical
partitions, compatible one-hot inputs expose any difference in weights.

Enumerate legal RHS axis assignments bottom-up, recording their
representations and witnesses. Finite local choices ensure termination;
induction on the tree ensures every legal assignment is represented.
Exactness therefore makes root matching sound and complete.
Separated addition follows componentwise from
Appendix~\ref{sec:addition-proof}.
\end{proof}

For fixed unary scaling alone, search is unnecessary. Removing scaling
nodes gives $E=c_E\bar E$, where $c_E$ is their product. The same one-hot
argument gives equivalence exactly when $c_L=c_R$ and
$N(\bar L)=N(\bar R)$. Thus compare coefficients and apply the original
direct constructor to the remaining trees, componentwise.
The general search may be combinatorial; this theorem does not establish
completeness of TEPL's implemented solver.
